\section{Related Work}
\label{sec:related_work}

This section reviews the core literature establishing our research context, organized into five thematic domains: YOLO-based detection evolution, visual state space models, hybrid CNN--Mamba architectures, edge deployment methodologies, and warehouse perception benchmarks.

\subsection{YOLO-Based Object Detection Evolution}
\label{sec:rel_yolo}
Single-stage object detectors originated with the unified regression formulation of YOLO \cite{redmon2016yolo}, which mapped image pixels directly to spatial bounding box coordinates and class conditional probabilities. Over the past decade, the YOLO family has evolved significantly \cite{sapkota2026ultralytics}. YOLOv5 established a modular anchor-based PyTorch architecture using Cross Stage Partial (CSP) networks and Path Aggregation Necks (PANet). YOLOv8 transitioned to an anchor-free design with a decoupled head, Distribution Focal Loss (DFL), and Task-Aligned Assigner (TAL) for dynamic label assignment. YOLO11 introduced C3k2 building blocks and cross-stage spatial attention (C2PSA) to optimize gradient pathways across model scales. Most recently, YOLO26 and YOLO27 have explored end-to-end NMS-free inference through dual-label assignment and scale-adaptive architectures \cite{sapkota2026ultralytics}. However, standard lightweight YOLO variants (e.g., YOLOv8n/s, YOLO11n/s) rely exclusively on local convolutions, which struggle to capture long-range contextual dependencies across visually complex, cluttered scenes without inflating latency.

\subsection{Mamba and State-Space Models for Visual Perception}
\label{sec:rel_mamba}
To overcome the quadratic computational scaling $O(N^2)$ of self-attention in Vision Transformers (ViTs), Selective State Space Models (SSMs) such as Mamba have emerged as an efficient sequence-modeling alternative operating with linear complexity $O(N)$ \cite{bao2026visionmamba}. While continuous state space models define 1D differential equations mapped through zero-order hold discretization, visual adaptations (e.g., Vision Mamba, VMamba, VSS-SpatioNet) translate 2D spatial feature maps into 1D token sequences using multi-directional scanning strategies (such as 4-way cross-scan SS2D) \cite{xiang2025vss, bao2026visionmamba}. These mechanisms enable global receptive fields with linear computational and memory scaling, making them promising building blocks for high-resolution visual processing.

\subsection{Hybrid CNN--Mamba Detection Architectures}
\label{sec:rel_hybrid}
Despite their theoretical strengths, pure visual SSMs face distinct challenges when directly applied to dense object detection: they lack the localized inductive bias of 2D convolutions and exhibit directional scanning blind spots. To resolve these limitations, recent works explore hybrid CNN--Mamba architectures:
\begin{itemize}
    \item \textbf{MambaNeXt-YOLO} \cite{lei2025mambanext}: Lei et~al.\ integrate ConvNeXt local feature extraction with an SS2D Mamba block and an adaptive ResGate fusion mechanism, complemented by a Multi-branch Asymmetric Fusion Pyramid Network (MAFPN). They report 85.1\% AP50 on PASCAL VOC (7.1M parameters, 22.4 GFLOPs) and 31.9 FPS on an NVIDIA Jetson Orin NX (TensorRT FP16).
    \item \textbf{Hybrid-YOLO} \cite{chen2024hybridyolo}: Chen and Er strategically embed Mamba blocks into the early stages of the YOLO backbone to capture global context, augmenting the design with lightweight spatial attention and Ghost-inspired efficient convolutions (+1.6 AP on COCO nano, 69.5 AP on underwater DUO).
    \item \textbf{Mamba-YOLO} \cite{wang2025mambayolo}: Wang et~al.\ introduce Regional Attention with Gated Enhancement in YOLO, proving that localized SSM branches effectively capture long-range spatial context without quadratic attention penalties.
    \item \textbf{UAVDet} \cite{yang2026uavdet}: Yang et~al.\ introduce the Cross Stage Partial Mamba (CSPMB) block and a Tiny-focused Feature Pyramid Network (TFPN), demonstrating that combining CSP structures with Mamba yields a +5.0\% small-object AP improvement while reducing parameters by 84.9\% on aerial imagery.
    \item \textbf{Xray-YOLO-Mamba} \cite{li2024xraymamba}: Li et~al.\ demonstrate that 2D selective state-space processing excels at resolving severe overlap and severe object occlusion where traditional convolutional attention filters fail.
\end{itemize}
As synthesized by Vo et~al.\ in their systematic literature review \cite{vo2026selectivereview}, visual state-space models hold enormous promise for edge robotics and real-time transportation detection, yet existing works predominantly target high-altitude aerial imagery or server-grade GPUs. How selective scans behave when surgically substituted into modern decoupled YOLO11 architectures under realistic multi-class warehouse occlusion remains unexamined.

\subsection{Lightweight Object Detection and Edge Deployment}
\label{sec:rel_edge}
Deploying deep learning detectors to edge platforms requires balancing algorithmic complexity with physical hardware constraints \cite{mittal2024lightweightsurvey, capra2022edgehardware}. As surveyed by Mittal \cite{mittal2024lightweightsurvey}, parameter counts and FLOPs alone are insufficient predictors of on-device latency: Memory Access Cost (MAC) and operator fragmentation (launching numerous small, non-fused kernels) frequently dominate runtime. In a comparative evaluation of runtime frameworks across edge platforms, Ratul et~al.\ \cite{ratul2024accelerating} proved that TensorRT and ONNX Runtime significantly outperform standard PyTorch by fusing operators and streamlining memory transfers. Crucially, Szegedi et~al.\ \cite{szegedi2026jetson} benchmarked modern object detectors on the NVIDIA Jetson Orin Nano, finding that FP16 reduced-precision execution in TensorRT cuts inference latency by up to $65\%$ over standard FP32 graphs while maintaining strict physical RSS bounds. Recent edge-tailored detectors such as LEAF-YOLO \cite{quang2024leafyolo} and Edge-YOLO \cite{li2023edgeyolo} further demonstrate that operator graph simplicity is vital for real-time edge execution.

\subsection{Warehouse Perception and Industrial Benchmarks}
\label{sec:rel_warehouse}
Industrial automation and logistics present severe visual perception challenges, including extreme object scale variation (from small packages to massive forklifts and pallet racks), dense occlusions, and variable lighting \cite{tan2024multiscale}. Investigating specialized warehouse environments, Yang et~al.\ \cite{yang2024freezerwarehouse} highlighted that high-density stacking and heavy visual overlap severely degrade conventional convolutional feature maps, necessitating focused attention or selective spatial context. Furthermore, large-scale benchmarks such as NVIDIA PhysicalAI-SmartSpaces \cite{wang2026smartspaces} provide vital multi-camera tracking and 3D perception environments across diverse warehouse facilities. Achieving real-time inference on edge devices in such environments requires detectors that capture room-scale aisle context while operating strictly within constrained memory envelopes.

\subsection{Summary and Research Gap Synthesis}
\label{sec:rel_gap}
While lightweight YOLO variants and visual state-space models have developed rapidly, prior hybrid CNN--Mamba detectors have been evaluated almost exclusively on generic natural images (PASCAL VOC, MS COCO) or remote sensing datasets (VisDrone, DOTA) \cite{lei2025mambanext, chen2024hybridyolo, yang2026uavdet}. No prior study has systematically investigated a lightweight hybrid YOLO-Mamba architecture explicitly optimized for multi-class 2D warehouse object detection under resource-constrained edge computing environments. This study addresses that defensible gap by proposing, evaluating, and benchmarking an edge-optimized hybrid detector against standard YOLO baselines.
